\chapter{振荡、时间筛选与适应}
双节点已经提供两条路线：正回路可以形成开关，负回路配上正自反馈则有二维振荡的入口。若没有正自边，前面的瞬时产生--一阶降解模型仍排除了周期轨道。第三个环节提供另一种时间结构：在负闭环中，它增加动态环节，让没有正自边的模型也能出现摆动失稳；在前馈中，它让同一输入沿快慢两条路径到达输出。下面分别讨论振荡、时间筛选与适应。
\begin{figure}[H]\centering
\begin{tikzpicture}
\node[gene] (a) at (0,1.4) {$X_1$};\node[gene] (b) at (3,1.4) {$X_2$};\node[gene] (c) at (1.5,-.5) {$X_3$};
\draw[rep] (a)--(b);\draw[rep] (b)--(c);\draw[rep] (c)--(a);
\node[gene] (u) at (6,1.4) {$u$};\node[gene] (y) at (9,1.4) {$y$};\node[gene] (z) at (7.5,-.5) {$z$};
\draw[act] (u)--(y);\draw[act] (u)--node[left]{慢支路}(z);\draw[rep] (z)--(y);
\end{tikzpicture}
\caption{左：三次抑制构成负反馈环。右：不一致前馈回路，直接激活与间接抑制在输出处竞争；输入 $u$ 在这里是外部信号。}
\end{figure}

\section{负反馈环：从线性失稳到稳定周期轨道}
三个相互抑制环节可写为（本节取相同的一阶降解率 $\gamma$）
\begin{equation}\dot x_1=f_1(x_3)-\gamma x_1,\quad\dot x_2=f_2(x_1)-\gamma x_2,\quad\dot x_3=f_3(x_2)-\gamma x_3.\end{equation}
反馈符号仍然是负，但稳定性已经不能照搬双节点的结论。要找出新的可能，先用稳态附近的斜率近似产生率曲线，计算三种基本偏差各自会增长还是衰减。记 $a_i=f_i'$ 为三个交叉斜率。三次抑制使 $a_1a_2a_3<0$，写成 $a_1a_2a_3=-k^3$、$k>0$，则
\[
\lambda I-J=\begin{pmatrix}\lambda+\gamma&0&-a_1\\-a_2&\lambda+\gamma&0\\0&-a_3&\lambda+\gamma\end{pmatrix},
\quad \det(\lambda I-J)=(\lambda+\gamma)^3-a_1a_2a_3.
\]
当乘积为 $-k^3$ 时，取 $-1$ 的三个立方根 $-1,\frac12\pm i\frac{\sqrt3}{2}$ 得到
\[
\lambda_1=-\gamma-k,\qquad\lambda_{2,3}=-\gamma+\frac k2\pm i\frac{\sqrt3k}{2}.
\]
$k<2\gamma$ 时，三种偏差的增长率都为负，会逐渐减小；$k>2\gamma$ 时，其中两种带摆动的偏差开始增长。门槛 $k=2\gamma$ 处，摆动的角频率为 $\sqrt3\gamma$，对应周期 $2\pi/(\sqrt3\gamma)$。
对对称 Hill 环 $f(s)=\alpha/(1+s^n)$，唯一对称稳态满足 $\alpha=\gamma s(1+s^n)$，因此
\[
\frac k\gamma=\frac{ns^n}{1+s^n}.
\]
这里也不会漏掉另外的不对称稳态：当三个调控函数相同且都严格下降时，沿环走过三次调控，浓度增加最终对应的是反馈输入下降。这条反馈关系与 $x$ 本身只会相交一次，而三个浓度相等的解已经满足关系，因此它就是整个环的唯一稳态。
当 $n\leq2$ 时不可能达到有限参数的阈值；当 $n>2$ 时
\[
s_c=\left(\frac2{n-2}\right)^{1/n},\qquad
\alpha_c=\gamma\frac n{n-2}\left(\frac2{n-2}\right)^{1/n}.
\]
在 $n=3,\gamma=1$ 下 $\alpha_c=3\sqrt[3]2$，演示的 $\alpha=20$ 位于稳态失稳一侧。
随 $\alpha$ 增加，$\dd s/\dd\alpha=1/[\gamma(1+(n+1)s^n)]>0$，且 $\dd k/\dd s>0$；所以随参数增大，这两种摆动方式确实会从衰减转为增长。这对复特征值的实部从负变正，数学上称为\emph{穿过虚轴}；“虚轴”就是复数平面中实部为零的位置。因为这里增长率随参数变化确实不为零，这次穿越是\emph{横截}的，也就是实际穿过去，而不是碰到边界后折返。这是 Hopf 分岔的候选点：稳态附近原本衰减的摆动可能变成周期振荡。

仅有复特征值穿过虚轴，还没有证明振荡会稳定下来。下面把 $n=3$ 的非线性一步也算完。取临界点 $s_c=\sqrt[3]2$、$\alpha_c=3\gamma s_c$，平移稳态后，系统在每个节点上的二、三阶项都是前一个节点偏差的平方与立方。记
\[
\beta_2=f''(s_c)=2\gamma s_c^2,\qquad
\beta_3=f'''(s_c)=-2\gamma s_c,\qquad
\omega_0=\sqrt3\gamma.
\]
这些系数由 $f(x)=\alpha_c/(1+x^3)$ 直接求导再代入 $s_c^3=2$ 得到。此处的 Hill 函数在正稳态附近光滑，具备展开所需的导数。

计算的目标是\emph{第一 Lyapunov 系数} $\ell_1$：它决定微小摆动增长后，三阶非线性会抑制振幅还是继续放大它。为了把这一计算写得可复核，令 $P(v_1,v_2,v_3)=(v_3,v_1,v_2)$，临界点的线性矩阵为 $A=-\gamma I-2\gamma P$。再定义逐分量相乘的运算 $\odot$，二、三阶项写成
\[
B(v,w)=\beta_2(Pv)\odot(Pw),\qquad
C(v,w,z)=\beta_3(Pv)\odot(Pw)\odot(Pz).
\]
这只是把三次相同的节点计算合并成一个记号；Taylor 展开中的因子仍为 $\frac12B$ 与 $\frac16C$。取
\[
\zeta=e^{-2\pi i/3},\quad q=(1,\bar\zeta,\zeta)^{\mathsf T},\quad
p=q/3,\qquad Aq=i\omega_0q,
\quad \langle p,q\rangle=1,
\]
其中 $\langle p,v\rangle=\bar p^{\mathsf T}v$，复数用来同时记录摆动的幅度和相位。标准 Hopf 系数公式（参考文献~6）给出
\begin{align*}
\ell_1=\frac1{2\omega_0}\operatorname{Re}\big\langle p,
&C(q,q,\bar q)-2B(q,A^{-1}B(q,\bar q))\\
&+B(\bar q,(2i\omega_0I-A)^{-1}B(q,q))\big\rangle.
\end{align*}
它综合了直接三阶作用，以及二阶作用经过其他恢复方向再传回摆动的影响。临界点另一个特征值为 $-3\gamma$，且 $2i\omega_0$ 不是特征值，所以式中的两个逆矩阵都有定义。

利用 $Pq=\zeta q$、$q\odot\bar q=\bm1$ 和 $q\odot q=\bar q$，三项分别可以直接算为
\[
\begin{split}
C(q,q,\bar q)&=\beta_3\zeta q,\\
B(q,A^{-1}B(q,\bar q))&=-\frac{\beta_2^2\zeta}{3\gamma}q,\\
B(\bar q,(2i\omega_0I-A)^{-1}B(q,q))&=\frac{\beta_2^2}{3i\omega_0}q.
\end{split}
\]
最后一项没有实部，$\operatorname{Re}\zeta=-1/2$，于是
\[
\ell_1=-\frac1{4\omega_0}
\left(\beta_3+\frac{2\beta_2^2}{3\gamma}\right)
=-\frac{5s_c}{6\sqrt3}<0.
\]
$\ell_1$ 的数值随特征向量的幅度约定改变，但其符号不变；上式对应这里明确给出的 $q,p$。负号说明非线性会限制微小振幅的增长。

现在可以应用\emph{Hopf 分岔定理}：系统光滑，只有一对简单纯虚特征值，其余特征值实部为负；前面已证明参数穿越是横截的，这里又算出 $\ell_1<0$。因此，在 $\alpha$ 略高于 $\alpha_c$ 时，唯一稳态附近确实产生一族\emph{轨道渐近稳定的周期解}。这就是超临界 Hopf 分岔。轨道稳定表示扰动后会回到同一条闭合轨道，允许在周期中的相位稍有变化；不是要求每个时刻都回到原轨迹的同一个点。

用振幅 $r$ 描述这一步，适当缩放后的局部方程具有 $\dot r=\mu r+\ell_1r^3+\cdots$ 的形式，$\mu$ 随 $\alpha-\alpha_c$ 同号。$\mu>0$ 时非零稳态振幅满足 $r^2\simeq-\mu/\ell_1$，其恢复斜率约为 $-2\mu<0$；振幅从零连续长出，周期趋于 $2\pi/\omega_0$。这把“稳态不稳定”“周期轨道存在”和“周期轨道能吸引附近轨迹”三件事连接成完整论证。

这个证明是\emph{临界点附近}的结论。图~\ref{fig:rep} 的 $\alpha=20$ 远离阈值，仍是给定参数下的数值演示；不能仅凭局部定理宣称所有更大参数都保留同一条稳定周期轨道。非线性系数的矩阵代入也在 \texttt{code/bistability\_details.py} 中复核。

\begin{figure}[H]\centering\includegraphics[width=.78\linewidth]{figures/repressilator.pdf}\caption{三节点负反馈环的数值轨迹，生成脚本见 \texttt{code/make\_figures.py}。}\label{fig:rep}\end{figure}

\section{快慢反馈：用迟滞开关形成可调时钟}
设相邻两个峰的时间间隔为 $T$，频率为 $f=1/T$，一周期内的最大、最小浓度为 $x_{\max},x_{\min}$，振幅在本文记为
\[
A_{\rm osc}=x_{\max}-x_{\min}.
\]
比较振荡器时，至少要分开三个问题：它能否持续振荡，频率能调多少，以及调频时振幅改变多少。一个只能在很窄的参数区间里产生微弱摆动的系统，与一个可以调节周期、又保持明确高低状态的系统，虽然都能画出周期轨迹，功能上并不相同。

也要说明调的是哪个参数。如果把\emph{所有}速率统一乘以同一系数，等于把整个时间坐标一起缩放，当然可以改变周期而保留轨迹形状。更有意义的问题是，只改变一个合适的产生率或时间尺度，其他过程不作同步调整时，是否仍能稳定调频。稳健性则是另一个量：它考察某个采样范围内，多少组参数产生符合标准的振荡；数值取决于模型、采样规则和振荡判据。

回到前一章的 R8 两变量模型，可以从快慢分离理解另一种振荡机制：
\[
\dot x=f(x,y),\qquad\dot y=\varepsilon g(x,y),
\]
其中 $f$ 与 $g$ 就是前一章的 R8 产生--降解规则。$\varepsilon$ 较小时，$x$ 快，抑制变量 $y$ 慢。暂时固定 $y$，只让 $x$ 恢复，就能画出\emph{快变量稳态曲线} $f(x,y)=0$。如果某个 $y$ 区间内存在明显分开的低表达和高表达两个稳定分支，中间夹着不稳定分支，$x$ 会在慢变量推动下沿一个稳定分支走，到折叠点才跳到另一支。

这就产生\emph{迟滞}：从低态转向高态与从高态转向低态，发生在不同的 $y$ 阈值。$y$ 增大时，高表达态一直保持到其分支终止，再快速降到低态；$y$ 减小时，低态又保持到另一个阈值，再快速升高。慢变化与快跳跃轮流发生，称为\emph{弛豫振荡}。这里“弛豫”是这一类振荡的标准名称，不表示整个过程都以单个指数恢复。

设高分支上 $y$ 从 $y_-$ 增到 $y_+$，低分支上再从 $y_+$ 降到 $y_-$，快分支上的稳态值分别记为 $x_H(y),x_L(y)$。若快慢分离足够明显、快跳跃所占时间很短，那么由 $\dd t=\dd y/[\varepsilon g(x,y)]$ 得到周期的近似：
\[
T\simeq\frac1\varepsilon\left[
\int_{y_-}^{y_+}\frac{\dd y}{g(x_H(y),y)}+
\int_{y_+}^{y_-}\frac{\dd y}{g(x_L(y),y)}\right].
\]
高分支上 $g>0$、低分支上 $g<0$，所以两项时间都为正。这说明慢变量的速度主要控制等待多久，而分支位置与跳跃阈值主要控制浓度跨越的范围。只改变 $\varepsilon$ 时，快稳态曲线本身不变，因此在这一近似适用的区间内，有可能明显改变周期而只小幅改变振幅。

\begin{figure}[H]\centering
\includegraphics[width=\linewidth]{figures/positive_feedback_tuning.pdf}
\caption{本文 R8 模型的快慢机制演示。左：快变量稳态曲线的稳定分支、不稳定分支与数值轨迹；中：改变同一个慢时间尺度，周期不同而浓度跨度接近；右：频率与振幅相对参考参数的变化。固定 $a_0=0.4,b_0=4,d_0=2$，仅改变 $\varepsilon$。}\label{fig:positive_tuning}
\end{figure}

这个近似有明确限度。快慢分离减弱、参数靠近分支消失的位置，或反馈改变了快稳态曲线本身时，振幅与周期仍可能一起变化。数值脚本在所示参数范围内检查了后期重复峰、正浓度与频率变化；这些数值检查说明例子怎样工作，并不替代极限环存在性的证明。

\section{前馈与适应：瞬态峰、支路抵消和误差积分}
输入 $u$ 可以直接激活输出 $y$，又通过慢变量 $z$ 抑制 $y$：
\begin{equation}\tau_y\dot y=H(u,z)-y,\qquad\tau_z\dot z=G(u)-z.\end{equation}
在 $H_u>0,H_z<0,G'>0$ 的调控范围内，若 $z$ 比 $y$ 慢，上升阶跃可使 $y$ 先增加，随后 $z$ 增加抑制 $y$，形成脉冲。\emph{脉冲}与\emph{精确适应}不同：精确适应要求稳态输出 $y^*(u)$ 与持续输入的幅度无关，这对 $H$、$G$ 有额外条件；拓扑本身不保证这种性质。后面将把支路抵消与积分反馈作具体比较。

要把“回落”与“回到同一基线”区分开，先算恒定输入下的长期输出。设 $\tau_y,\tau_z>0$，$H,G$ 在所考察输入及状态范围内连续，固定 $u$ 时有
\[
z(t)=G(u)+(z_0-G(u))e^{-t/\tau_z},
\]
所以 $z\to G(u)$。此时输入输出支路 $H(u,z(t))\to H(u,G(u))$，按前面单向调控的方法把 $y$ 的方程积分，就得到
\[
y^*(u)=H(u,G(u)).
\]
精确适应要求输入改变后仍回到相同的输出，因此恰好要求这个复合函数在输入区间内为常数。如果 $H,G$ 为 $C^1$，且所讨论的输入覆盖一个完整区间 $I$，那么可以用\emph{链式法则}沿两条支路计算输入变化的影响。链式法则就是：间接支路的总变化率等于沿路各步变化率的乘积。由此得到
\[
H_u(u,G(u))+H_z(u,G(u))G'(u)=0\qquad(u\in I).
\]
这个式子揭示了适应所需的抵消：直接输入对输出的长期增益，与输入通过 $z$ 产生的长期增益大小相等、符号相反。快慢差异帮助形成瞬态峰，但只有这一稳态抵消才使基线与输入幅度无关。

可以选一个简单函数把这两个层面同时展示出来。取 $u>0,z>0$、$G(u)=u$、$H(u,z)=u/z$。每个恒定输入的稳态为 $(y^*,z^*)=(1,u)$，故确有精确适应。若从输入 $u_0$ 的稳态出发，在 $t=0$ 切换至 $u_1>u_0$，则
\[
z(t)=u_1+(u_0-u_1)e^{-t/\tau_z},\qquad
\dot y(0^+)=\frac{u_1/u_0-1}{\tau_y}>0.
\]
$y$ 起初升高，最终回到 $1$，所以至少存在一个高于基线的峰；$\tau_z\gg\tau_y$ 使脉冲更易分辨。$u/z$ 是理想化的比值调控，只在 $u,z>0$ 的工作域内使用。若换成 $H=u/(1+z)$，保持同一拓扑与 $G$，则 $y^*=u/(1+u)$ 依赖输入，精确适应消失。这给出了拓扑不充分的明确反例。
\begin{figure}[H]\centering\includegraphics[width=.85\linewidth]{figures/feedforward_adaptation.pdf}
\caption{相同的不一致前馈结构可以精确适应，也可以只产生瞬态峰而改变最终基线。输入从 $1$ 跳至 $2$，$\tau_y=0.25,\tau_z=2$；右图比较两种调控函数的稳态输入输出曲线。}\label{fig:adapt}
\end{figure}


前馈回路有两条从输入到输出的路径。若两条路径的累计符号相同，称为\emph{一致前馈}；相反则为\emph{不一致前馈}。前面的例子利用两条路径相反的长期作用形成适应，而一致前馈的一种用途是区分短脉冲与持续输入。

看一个全激活的三节点回路 $A\to B\to C$、$A\to C$。为单独看时间逻辑，暂时把输入节点 $A$ 的活动当成 $0$ 或 $1$，令 $\tau_B\dot B=A-B$，再以阈值 $0<\theta<1$ 判断 $B$ 是否足够高。如果输出产生采用 AND 逻辑，要求 $A$ 与 $B$ 同时高：上升输入到来后，$B(t)=1-e^{-t/\tau_B}$，只有等到
\[
t_{\rm on}=-\tau_B\ln(1-\theta)
\]
才开始产生输出；输入一撤去，AND 条件立即失效。所以比这个时间还短的输入可能被过滤掉。如果改成 OR 逻辑，输入到来时直接支路立即打开，撤去后则要等 $B$ 从高值降到阈值，关闭的等待时间为 $t_{\rm off}=\tau_B\ln(1/\theta)$。

这是理想阈值模型，真实的连续产生率会把“立即”与“等待”平滑成不同快慢的响应，输出自身的恢复时间也会再加一层滞后。重点是：同样的三条激活边，配上不同启动子逻辑，就可能分别延迟开启或延迟关闭。因此一致前馈并非“没有适应所以无趣”，它回答的是另一个时间功能问题。


为什么普通的负反馈不自动给出精确适应？因为一个通常的恢复变量会同时遗忘过去的偏差：产生与降解平衡时，允许输出停在一个小的残余误差上。要把这个误差消掉，调节变量需要在误差存在时继续改变。

下面用一个三变量模型说明积分反馈怎样实现这一作用。$A$ 接收输入，$C$ 是输出，$B$ 是调节变量；$C$ 激活 $B$，$B$ 抑制 $C$。令
\begin{align}
\dot A&=\rho\{u(1-A)-A\},\\
\dot B&=\kappa(C-C_0),\\
\dot C&=k\{A(1-C)-BC\},\qquad \rho,\kappa,k>0.\label{eq:buffer_demo}
\end{align}
$A,C$ 可以理解为活性比例；$B$ 是缓冲变量，在本例的工作范围内也保持正且小于 $1$。第一式使 $A^*=u/(1+u)$。关键在第二式：当 $C>C_0$，$B$ 继续积累并压低 $C$；当 $C<C_0$，$B$ 下降，抑制减弱；只有 $C=C_0$ 时，$B$ 才停止调整。因此
\[
C^*=C_0,\qquad B^*=\frac{1-C_0}{C_0}A^*.
\]
输入变大后，长期改变的是 $A^*$ 与 $B^*$，输出仍为 $C_0$。这不是两个支路恰巧抵消一次，而是反馈持续校正误差。

把第二式积分，得到
\[
B(t)=B(0)+\kappa\int_0^t[C(s)-C_0]\dd s.
\]
所以 $B$ 保存的是过去输出误差的累积，称为\emph{积分反馈}。这里的“缓冲”不只是响应慢，而是 $\dot B$ 在工作近似下不直接依赖 $B$ 自己。饱和的酶反应能提供近似恒定的转化速率，从而产生这类条件；它不会无条件地出现在所有负反馈回路中。式~\eqref{eq:buffer_demo}采用线性积分近似；若 $B$ 达到上下限，线性积分近似就失效，必须回到完整的转化反应。

这个例子的恢复也能算出来。因为 $A$ 不受下游影响，线性化后它的一根特征值为 $-\rho(1+u)$，而 $(B,C)$ 部分为
\[
J_{BC}=\begin{pmatrix}0&\kappa\\-kC_0&-k(A^*+B^*)\end{pmatrix}.
\]
迹为负，行列式 $\kappa kC_0>0$，故只要 $0<C_0<1$ 且 $B^*$ 在缓冲近似的内部工作范围，稳态就局部渐近稳定。输出基线固定与基线附近能恢复，是这里分别检查过的两件事。

若给缓冲变量加入“泄漏”，改成
\[
\dot B=\kappa(C-C_0)-\varepsilon_B(B-B_{\rm ref}),\qquad\varepsilon_B>0,
\]
稳态就允许 $C^*-C_0=(\varepsilon_B/\kappa)(B^*-B_{\rm ref})$。它仍是负反馈，也仍可能有瞬态峰，却通常留下随输入变化的基线偏差。缓冲变量慢并不足够，是否能持续积累误差才是区别。


前馈的思路不同：没有输出回传的校正，而是让一条支路产生与输入成比例的反向作用。沿用同样的输入节点 $A$，构造另一个简化模型：
\begin{align}
\dot A&=\rho\{u(1-A)-A\},\\
\dot B&=k_BA-\gamma_BB,\\
\dot C&=k_AA(1-C)-k_CB C,\qquad k_B,\gamma_B,k_A,k_C>0.
\end{align}
$A$ 直接激活 $C$，又通过 $B$ 抑制 $C$。令 $\eta=k_B/\gamma_B$，稳态下
\[
B^*=\eta A^*,\qquad
0=A^*\{k_A(1-C^*)-k_C\eta C^*\}.
\]
在 $A^*>0$ 的工作范围内，公共因子 $A^*$ 可以约去，得到
\[
C^*=\frac{k_A}{k_A+k_C\eta}.
\]
这就是比例节点的作用：输入增加多少，抑制支路长期也按固定比例增加多少，两条支路在输出方程中共享输入因子。$B$ 较慢时，直接支路先起作用，随后比例支路跟上，于是输出先偏离、再回到基线。

该模型是逐级驱动的：$A$ 先稳定，随后 $B$ 稳定，最后 $C$ 的恢复率为 $k_AA^*+k_CB^*>0$，所以在给定正输入及上述工作范围内三者都收敛。若把 $B$ 的产生换成 $k_BA/(K+A)$，则 $B^*/A^*=k_B/[\gamma_B(K+A^*)]$ 不再恒定，输出基线也通常随输入改变。可见不一致前馈的边符号还要配上合适的支路比例关系。

\begin{figure}[H]\centering
\begin{tikzpicture}
\node[gene] (a) at (0,1.3) {$A$};\node[gene] (c) at (3,1.3) {$C$};\node[gene] (b) at (3,-.4) {$B$};
\draw[act] (a)--(c);\draw[act] (c) to[bend left=22] (b);\draw[rep] (b) to[bend left=22] (c);
\node at (1.5,-1.15) {缓冲：$\dot B\propto C-C_0$};
\node[gene] (aa) at (6,1.3) {$A$};\node[gene] (cc) at (9,1.3) {$C$};\node[gene] (bb) at (7.5,-.4) {$B$};
\draw[act] (aa)--(cc);\draw[act] (aa)--(bb);\draw[rep] (bb)--(cc);
\node at (7.5,-1.15) {比例：$B^*\propto A^*$};
\end{tikzpicture}
\caption{两种适应机制的节点角色。左边由输出误差驱动缓冲变量持续校正；右边由输入驱动比例变量，长期抵消直接输入作用。两者都允许 $B$ 的稳态随输入改变。}
\end{figure}
\begin{figure}[H]\centering\includegraphics[width=\linewidth]{figures/adaptation_mechanisms.pdf}
\caption{本文简化模型的数值比较。左：缓冲反馈与比例前馈都先响应再恢复，带泄漏的缓冲反馈改变最终基线。中：调节变量 $B$ 保留输入变化。右：理想机制的稳态输出与输入无关，泄漏模型则不然。曲线按各自初始基线归一化。}\label{fig:adaptmech}
\end{figure}
本图的示例参数为 $\rho=3$、输入 $u:0.5\to0.8$。缓冲模型取 $C_0=0.4,\kappa=0.3,k=2$，泄漏模型另取 $\varepsilon_B=0.08$；比例模型取 $k_B=\gamma_B=0.25,k_A=k_C=2$。这些参数用于展示基线恢复与泄漏反例。完整方程与检查见 \texttt{code/extended\_examples.py}。

同一种连接也可能通过不同变量保存误差。例如在 $B>0$ 的工作域内，若调节规律为 $\dot B=\nu B(C_0-C)$，则
\[
\frac{\dd\log B}{\dd t}=\nu(C_0-C).
\]
累积误差的变量现在是 $\log B$；内部稳态仍要求 $C=C_0$。这里的调节规律涉及随状态改变的产生或损失，必须另外检查正性、工作范围和稳定性，不能直接套用固定一阶降解下的拓扑分类。

精确适应与参数稳健性也不同。前者要求同一组参数下，稳态输出不随持续输入改变；后者问这种性质能否在一个参数范围内保持。积分反馈的基线由 $C_0$ 指定，比例前馈的基线则由速率比指定：改变这些设定参数可以改变基线，却不必破坏对输入的适应。因此分析功能时要同时指定节点角色、速率机制、参数和输入方式。
